\cotefile{10}
\scanpage{139}
Motifs à coefficients sur un corps de nbs.

Soit d'abord $\unsure{\mathcal C|C|\underline C}$ une $\otimes$-catégorie sur un corps $K$. Si $K'$ est
une extension \inserted{\uline{finie}} de $K$, on désigne par $\mathcal C^{K'}$ la catégorie des
objets de $\mathcal C$ sur lesquels $K'$ opère \unsure{linéairement}, \unsure{i.e.}
des $X\in\operatorname{Ob}\mathcal C$ munis d'un hom de $K$-algèbres
$K'\longrightarrow\operatorname{End}(X)$. C'est encore une $\otimes$-catégorie, en définissant
$X\otimes_{K'}Y$ comme le quotient de $X\otimes_KY$ par les
$\lambda\otimes\mathrm{id}_Y-\mathrm{id}_X\otimes\lambda$ $(\lambda\in K')$, et $\underline{\mathrm{Hom}}_{K'}(X,Y)$ comme le
noyau \unsure{dans|des|de} $\underline{\mathrm{Hom}}_K(X,Y)$ des $\underline{\mathrm{Hom}}(\lambda_X,\mathrm{id}_Y)-\underline{\mathrm{Hom}}(\mathrm{id}_X,\lambda_Y)$.

Quand $\mathcal C$ est défini par un schéma en groupes
affine $G$ sur $K$\inserted{, grâce au foncteur fibre $F$ sur $\mathcal C$}, alors $\mathcal C^{K'}$ est défini par
le groupe affine $G_{K'}=G\otimes_K K'$, grâce au foncteur
fibre \inserted{sur $\mathcal C^{K'}$} $F_{K'}$ déduit trivialement de $F$.

On a un $\otimes$-foncteur naturel $\mathcal C\longrightarrow\mathcal C^{K'}$, \unsure{noté|notant}
$X\leadsto X\otimes_K K'$, et dont le composé avec le foncteur
``oubli'' est bien $X\otimes_K\underset{\sim}{K'}$ [où $\underset{\sim}{K'}$ est \deleted{\illegible{1}} \unsure{le'objet|l'objet}
de $\mathcal C$ défini par le \uline{vectoriel} $K'$ sur $K$, -- puisqu'on
a un $\otimes$-foncteur évident Vectoriels \inserted{de dim finie} sur $K\longrightarrow\mathcal C$].

Supposons $\mathcal C$ la catégorie des \inserted{(iso)}motifs ($1/2$ simples) sur un
corps $k$, \unsure{donc} $K=\mathrm Q$, $K'$ un corps de nbs.
Soit $E$ un motif à coefficients dans $K'$, donc
un motif ordinaire, avec $K'\longrightarrow\operatorname{End}_{\mathrm Q}(E_0)$, $E_0$ le $\mathrm Q$-motif \unsure{ss-jacent}.
Alors $\unsure{\check E|\widetilde E}$ \inserted{dans $\mathcal C^{K'}$} est défini par $\widetilde E$ dans $\mathcal C$, avec l'hom
% archive sheet number: 138
\scanpage{140}
\unsure{transposé|donné par} $K'\xrightarrow{\varphi}\operatorname{End}_{\mathrm Q}(E_0)\xrightarrow{u\leadsto{}^tu}\operatorname{End}_{\mathrm Q}(\check E_0)$. Donnons
nous une polarisation sur $E_0$ (le $\mathrm Q$-motif \deleted{\unsure{ss-jacent}} à $E$).
Alors on trouve un isom $\check E_0\simeq E_0(\unsure{p|\rho})$
($E_0$ de poids $p$), donc \deleted{$K'\longrightarrow\operatorname{End}_{\mathrm Q}(\check E_0)$}
$\operatorname{End}_{\mathrm Q}(\check E_0)\simeq\operatorname{End}(E_0(p))\simeq\operatorname{End}(E_0)$, et on a donc défini
une \uline{involution} \inserted{$u\leadsto u'$} sur $\operatorname{End}_{\unsure{\mathrm Q}}(E_0)$. Le motif dual $(\deleted{\unsure{E_\varphi}}\!E_\varphi)^\vee$
\inserted{de $E_\varphi$ (défini par $\varphi$)} est défini par l'hom $\lambda\leadsto\varphi(\lambda)'\underset{\unsure{0|\circ}}{=}\varphi'(\lambda)$ comme
\[
\boxed{(E_\varphi)^\vee=(E_{\varphi'})(p)}
\]
On n'a pas en général un isom $(E_\varphi)^\vee\simeq E_\varphi(p)$ !

Prenons l'exemple des $K'$-motifs de niveau zéro,
poids $0$. La catégorie est équivalente à la catégorie
des vectoriels de dim finie sur $K'$, \deleted{\illegible{1}}
\deleted{\unsure{en}} à opération \inserted{continue} de $\pi=\operatorname{Gal}(\bar k/k)$. Le dual
correspond à l'opération contragrédiente de $\pi$. Les
$K'$-motifs de rang $1$ \inserted{à isom près} correspondent aux
représentations \deleted{de} \inserted{continues} $\pi\xrightarrow{\chi}K'{}^*$, et on a
$(L_\chi)^\vee=L_{\check\chi}$, où $\check\chi(g)=\chi(g^{-1})$. Donc
on a $L_\chi{}^\vee\simeq L_\chi$ ssi $\chi=\check\chi$ i.e. $\chi$ est
à valeurs dans $\pm1$ !
% archive sheet number: 139
\scanpage{141}
\deleted{\uline{\unsure{Fascicule|Formalisme} \unsure{tensoriel|finie}}}

Structures supplémentaires sur des vectoriels, et algèbres de Hopf top.

Anneaux des représentations de certains groupes Algébriques.
% archive sheet number: 140
\scanpage{143}
Considérons la catégorie des espaces vectoriels \inserted{de dim finie} sur un corps $k$
(pour fixer les idées). Pour tout groupe fini $G$, soit
${\unsure{\mathcal A|A}}^G$ la catégorie des espaces vectoriels de dim finie où
$G$ opère. C'est \deleted{aussi} une catégorie avec produit
tensoriel \inserted{(\uline{exact} svt.)} et involution \inserted{(et op. $\lambda$)} (et $\mathcal A^{\{e\}}$ est la catégorie de départ).

\deleted{On a de plus un bifoncteur $\mathcal A^G\times\mathcal A^{G'}\longrightarrow\mathcal A^G$}

Cette catégorie dépend fonctoriellement de $G$ : si on a
un hom $u:G\longrightarrow G'$, on a un foncteur \uline{contravariant}
\[
\deleted{\unsure{\mathrm{Res}}}\,\mathcal A^u:\mathcal A^{G'}\longrightarrow\mathcal A^G \tag*{(1)}
\]
% slanted
\margin{$\mathcal A^u$ est un ``isomorphisme dans'' si $u$ est surjectif}
\[
\deleted{\unsure{(2)}}\,\mathcal A^{\unsure{1_G|\mathrm{id}_G}}=\deleted{\unsure{\mathcal A}}\,1_{\mathcal A^G}\ ;\quad \mathcal A^{uv}=\mathcal A^v\mathcal A^u \tag*{$(1')$.}
\]
\deleted{(le foncteur est contravariant)}. Cela permet de définir de façon purement
formelle le produit tensoriel externe
\[
\mathcal A^G\times\mathcal A^{G'}\longrightarrow\mathcal A^{G\times G'} \tag*{(2)}
\]
en posant, si $p_1,p_2$ sont les projections de $G\times G'$ sur ses
deux facteurs
\[
M\boxtimes M'=\deleted{\unsure{p_1^*(M)\otimes p_2^*(M')}}\,\mathcal A^{p_1}(M_1)\otimes\mathcal A^{p_2}(M'_2) \tag*{$(2')$}
\]
De \unsure{même} on définit la somme directe extérieure $\boxplus$
\deleted{(on peut donc dire}
% slanted
\margin{dans un sens bien compliqué à préciser !!}

$\boxtimes$ est un bifoncteur additif, \inserted{``commutatif'', ``\uline{associatif}''.} Sa connaissance
(indépendamment des structures multiplicatives dans les
$\mathcal A^G$) permet de reconstruire ces dernières, par
considération de $\Delta_G:G\longrightarrow G\times G$ : on a
\[
M\otimes M'=\deleted{\unsure{\Delta_G^*}}\,\mathcal A^{\Delta_G}(M\boxtimes M') \tag*{(3)}
\]
Enfin, on a aussi une notion de \uline{représentation}
\uline{induite}, qui pour un groupe \deleted{$G$ et un ss-groupe} \inserted{hom \uline{injectif}}
$u:G\longrightarrow G'$, définit un foncteur \uline{covariant} additif
% archive sheet number: 142
% printed form (paper watermark, vertical, cut off at both ends)
\unsure{YPEWRITER} \tikz[baseline=-.4ex,scale=.15]{\foreach \a in {0,45,...,315}{\draw (\a:.52) circle (.22);} \draw (0,0) circle (.21);} ARGENTINE \tikz[baseline=-.2ex,scale=.17]{\draw (0,0)--(1.6,0)--(1.8,.33)--(.5,.33)--(.32,.13)--(.92,.13);} \unsure{TYPEW}
\scanpage{144}
\[
I^u:{\unsure{\mathcal A|A}}^G\longrightarrow\mathcal A^{G'} \tag*{(4)}
\]
On a aussi \deleted{On a aussi \unsure{$(2')$}}
\[
I^{1_G}=\deleted{\unsure{I}}\,1_{\mathcal A^G},\quad I^{uv}=I^uI^v \tag*{$(4')$}
\]
Les relations entre ces deux sortes de \deleted{opérations} \inserted{foncteurs} $\mathcal A^u,I^u$
et les \deleted{foncteurs} opérations $\otimes$ sont
\[
\left\{\begin{aligned}
\mathcal A^u(I^u(M))&=M^{[G':G]}\\
\deleted{\unsure{I^u(\mathcal A^u(M'))=M'{}^{[G':G]}}}&
\end{aligned}\right.\quad(M\in\mathcal A^G,\ G\subset G') \tag*{(5)}
\]
\deleted{\unsure{(6)}} On a d'ailleurs un hom fonctoriel
\deleted{\unsure{(6)} $M\longrightarrow$}
\[
M\longrightarrow\mathcal A^u I^u(M) \tag*{(6)}
\]
\[
M'\longrightarrow I^u\mathcal A^u(M') \tag*{(7)}
\]
\inserted{\unsure{$G$|2} hom. naturels !}

\margin{(8) Si $u$ est un isom, alors $I^u\unsure{=|\simeq}\mathcal A^{u^{-1}}$, donc $I^u$ et $\mathcal A^u$ sont des foncteurs inverses l'un de l'autre.}

\deleted{\unsure{Définissons|Définition}} (9) Soit $s$ la symétrie $G_1\times G_2\longrightarrow G_2\times G_1$. Alors $I^s(M\boxtimes M')=M'\boxtimes M$
% margin note written sideways
\margin{(10) Si $g\in G$, et $\sigma_g$ est l'automorphisme \uline{intérieur} défini par $g$, on a pour $M\in\mathcal A^G$ $M\simeq I^{\sigma_g}(M)\simeq\mathcal A^{\sigma_g}(M)$ L'isomorphisme \unsure{it.|itou|étant} est canonique et défini par $g$ opérant sur le vectoriel sous-jacent à $M$.}
% margin note written sideways
\margin{\unsure{formuler la propriété de cet isomorphisme}.}

Soit, pour des entiers $\geq1$, $p_1\ldots p_n$ de somme $p$
\[
u_{p_1,\ldots,p_n}:\unsure{\mathfrak S|\mathcal S}_{p_1}\times\ldots\times\mathfrak S_{p_n}\longrightarrow\mathfrak S_p \tag*{(\unsure{11|10})}
\]
l'hom (injectif) bien connu. Posons
pour $M\in\mathcal A^{\mathfrak S_p}$, $M'\in\mathcal A^{\mathfrak S_q}$
\[
M*M'=I^{u_{p,q}}(M\boxtimes M')\deleted{\unsure{\in\mathcal A^{\mathfrak S_{p+q}}}} \tag*{(\unsure{12|11}).}
\]
On a un isomorphisme canonique
\[
M*M'\simeq M'*M \tag*{(\unsure{13|12})}
\]
\deleted{canonique}, \deleted{\unsure{mais} à vrai dire \unsure{n'étant un tel} isomorphisme \unsure{qui est}} \inserted{défini par}
$\alpha$ : il existe $\alpha$ \deleted{des} unique de $\mathfrak S_{p+q}$
tel que $\alpha u_{p,q}\alpha^{-1}=u_{q,p}s_{p,q}$
-- où $s_{p,q}$ est la symétrie $\mathfrak S_p\times\mathfrak S_q\longrightarrow\mathfrak S_q\times\mathfrak S_p$.
% slanted
\margin{$*$ Noter aussi $\mathcal A^{u_{p,q}}(M*M')=\deleted{\unsure{M\boxtimes M'}}(M\boxtimes M')^{\frac{(p+q)!}{p!\,q!}}$ ce qui relie le $\boxtimes$ et le $*$, on peut \unsure{écrire} \deleted{$(M\otimes M')^{\unsure{\frac{(2p)!}{(p!)^2}}}$} $(M\odot M')^{\unsure{(2p)!/(p!)^2}}$ est défini par}

De façon analogue, on a un isomorphisme
canonique
\[
(M*M')*M''\simeq M*(M'*M'')\simeq I^{u_{p,q,r}}(M\boxtimes M'\boxtimes M'') \tag*{(14)}
\]
% archive sheet number: 143
% printed form (paper watermark visible in reverse)
\tikz[baseline=-.4ex,scale=.15]{\foreach \a in {0,45,...,315}{\draw (\a:.52) circle (.22);} \draw (0,0) circle (.21);} ARGENTINE \tikz[baseline=-.2ex,scale=.17]{\draw (0,0)--(1.6,0)--(1.8,.33)--(.5,.33)--(.32,.13)--(.92,.13);} TYPEWRITER \tikz[baseline=-.4ex,scale=.15]{\foreach \a in {0,45,...,315}{\draw (\a:.52) circle (.22);} \draw (0,0) circle (.21);} \unsure{A|AR}
\scanpage{145}
\deleted{\unsure{Dans le domaine de la} \illegible{1}}

Soit $M\in{\unsure{\mathcal A|A}}^G$, alors pour $\forall$ entier $n$, $\bigotimes^n M$ est un
élément de $\mathcal A^{G\times\unsure{\mathfrak S|\mathcal S}_n}=(\mathcal A^G)^{\mathfrak S_n}$. En \deleted{\unsure{tout}} \unsure{plus bref|par abr.},
on notera cet élément par $T^n(M)$. Ainsi $T^n$ est
un foncteur (non plus additif !)
% slanted
\margin{Expliquer la \unsure{variation|opérations} : on considère $\boxtimes^n M\in\mathcal A^{G^n}$, où $\mathfrak S_n$ opère dans $G^n$ et dans $\boxtimes^n M$\ldots}
% slanted
\margin{Dire que si un \unsure{gpe|groupe} fini $G'$ opère dans un élément de $\mathcal A^G$, on obtient un élément de $\mathcal A^{\unsure{G\rtimes G'|G\times G'}}$}
\[
T^n:\mathcal A^G\longrightarrow(\mathcal A^G)^{\mathfrak S_n} \tag*{(15)}
\]
Sa propriété essentielle est
\[
\boxed{T^n(M+M')\simeq\sum_{p+q=n}T^p(M)*T^q(M')} \tag*{(16)}
\]
% slanted
\margin{et propriété analogue pour les \uline{\unsure{diviseurs|divisés}} ?}
De plus, $T^n$ commute aux $\mathcal A^u$ : Soit $u:G\longrightarrow G'$,
alors le diagramme de foncteurs
\diagram{\begin{tikzcd}[ampersand replacement=\&,column sep=huge]
\text{(17)} \& \mathcal A^{G'} \arrow[r,"T^n_{G'}"] \arrow[d,"\mathcal A^u"'] \& (\mathcal A^{G'})^{\mathfrak S_n} \arrow[d,"\mathcal A^{u\times 1_{\mathfrak S_n}}"]\\
\&\mathcal A^G \arrow[r,"T^n_G"] \&(\mathcal A^G)^{\mathfrak S_n}
\end{tikzcd}}
est commutatif (à préciser les compatibilités
\deleted{\unsure{des}} éventuelles d'hom. fonctoriels).

Considérons
\[
v_{p,q}:\mathfrak S_p\times\mathfrak S_q\longrightarrow\mathfrak S_{pq} \tag*{(18)}
\]
défini par $v_{p,q}(\sigma,\tau)(i)=\alpha_{p,q}(\sigma)\,u_{\underbrace{\tau,\ldots,\tau}_{p}}\,\alpha_{p,q}(\sigma)^{-1}$
où $\alpha_{p,q}\deleted{(\sigma)}:\mathfrak S_p\longrightarrow\mathfrak S_{pq}$ est défini par
\[
\left\{\begin{aligned}
&\alpha_{p,q}(\sigma)\text{ est croissante dans les intervalles } ]iq,i(q+1)]\\
&\text{et l'applique sur l'intervalle}\\
&]\sigma(i)q,\sigma(i)(q+1)]
\end{aligned}\right. \tag*{(19)}
\]

Sous ces conditions, on a
\[
T^m(T^n(M))\simeq\deleted{\unsure{v_{m,n}}}\,\mathcal A^{v_{m,n}}(T^{mn}(M)) \tag*{(20)}
\]
(les deux mbres dans $\mathcal A^{G\times\mathfrak S_m\times\mathfrak S_n}$)
% archive sheet number: 144
% printed form (vertical paper watermark)
\unsure{NTINE} \tikz[baseline=-.2ex,scale=.17]{\draw (0,0)--(1.6,0)--(1.8,.33)--(.5,.33)--(.32,.13)--(.92,.13);} TYPEWRITER \tikz[baseline=-.4ex,scale=.15]{\foreach \a in {0,45,...,315}{\draw (\a:.52) circle (.22);} \draw (0,0) circle (.21);} \unsure{ARGENTINE|ARGENTIN}
\scanpage{146}
On a aussi
\[
T^n(M\otimes M')\simeq T^{\unsure{m|n}}(M)\otimes T^n(M')\quad\in{\unsure{\mathcal A|A}}^{G\times\unsure{\mathfrak S|\mathcal S}_n} \tag*{(21)}
\]
% archive sheet number: 145
% printed form (vertical paper watermark, reverse face)
\unsure{NTINE} \tikz[baseline=-.2ex,scale=.17]{\draw (0,0)--(1.6,0)--(1.8,.33)--(.5,.33)--(.32,.13)--(.92,.13);} TYPEWRITER \tikz[baseline=-.4ex,scale=.15]{\foreach \a in {0,45,...,315}{\draw (\a:.52) circle (.22);} \draw (0,0) circle (.21);} \unsure{ARGENTINE|ARGENTIN}
\scanpage{147}
De plus, le foncteur $M*M'$ est bi-additif en chaque
argument. Il y a un élément unité pour $*$ :
$K\in{\unsure{\mathcal A|A}}^{\unsure{\mathfrak S|\mathcal S}_0}$, \deleted{\illegible{1}} [où par définition $\mathfrak S_0=\{e\}$, et $u_{0,p,\ldots}=u_{p,\ldots}$]

Passons aux groupes de classes (pour \unsure{le splittage|les splittings},
pour fixer les idées), posons
\[
\deleted{\unsure{K(G)}}\quad A(G)=K(\mathcal A^G) \tag*{(22)}
\]
$A(G)$ est un \deleted{\unsure{faisceau}} \uline{anneau de Chern}
\uline{\unsure{involutif|commutatif}} \uline{augmenté} [de plus, il est de dimension
finie, et même de rang fini sur les entiers, mais
c'est fortuit], $A(G)$ est un foncteur
\deleted{\uline{covariant}} contravariant en $G$, $u:G\longrightarrow G'$ définit
\[
A(u):A(G')\longrightarrow A(G) \tag*{(\unsure{23})}
\]
% slanted
\margin{$A(u)$ est injectif si $u$ est surjectif (grâce au th. de Jordan-Hölder)}
et $A(1_G)=1_{A(G)}$, $A(uv)=A(v)A(u)$. Cela permet de
définir le produit extérieur $\alpha\in A(G)$, $\beta\in A(G')$
\[
\alpha\square\beta=\deleted{\unsure{\alpha}}\ A(p_1)(\alpha).\,A(p_2)(\beta)\in A(G\times G') \tag*{(\unsure{24})}
\]
D'ailleurs, la seule connaissance de la structure de
groupe sur les $A(G)$, et de l'opération externe
$\alpha\square\beta$, redonne les structures d'anneau \deleted{\illegible{1}} des $A(G)$.
[Mais pas leur structure de $\lambda$-anneau pour l'instant).

Il y a aussi les foncteurs \inserted{\uline{additifs}}, pour $u:G\longrightarrow G'$
\unsure{injectifs}
\[
I(u):A(G)\longrightarrow A(G') \tag*{(25)}
\]
(mais ce foncteur n'\uline{est pas multiplicatif})

On a
\[
\begin{aligned}
A(u)I(u)&=\text{multiplication par }[G':G]\\
I(u)A(u)&=\text{multiplication par }[G':G]
\end{aligned}\tag*{(26)}
\]
% archive sheet number: 146
% printed form (vertical paper watermark)
\unsure{GENTINE} \tikz[baseline=-.2ex,scale=.17]{\draw (0,0)--(1.6,0)--(1.8,.33)--(.5,.33)--(.32,.13)--(.92,.13);} TYPEWRITER \tikz[baseline=-.4ex,scale=.15]{\foreach \a in {0,45,...,315}{\draw (\a:.52) circle (.22);} \draw (0,0) circle (.21);} \unsure{ARGENT}
\scanpage{148}
Ceci implique que si $u$ est un isomorphisme, alors
$A(u)$ et $I(u)$ sont des isomorphismes inverses l'un
de l'autre. Enfin
\[
A(\sigma_g)=I(\sigma_g)=\text{identité dans } A(G)\quad(\text{si }g\in G). \tag*{(27)}
\]

Ces propriétés suffisent pour mettre sur la somme
directe $A$ des $A^n$ :
\[
\begin{aligned}
A^n&=A(\unsure{\mathfrak S|\mathcal S}_n)\quad(n\geq1)&\big|&& A^*&=\coprod_{n=0}^{\infty}A_{\unsure{n}}\\
A^0&=\mathbb Z
\end{aligned}\tag*{(28)}
\]
une structure \deleted{\unsure{de}} anneau gradué commutatif.
\inserted{(et de même sur les $A(G)^n=A^G(\mathfrak S_n)$, $A(G)^*=\coprod_{0}^{\infty}A(G)^n$}

D'ailleurs, on a, pour $\alpha\in A^p$, $\beta\in A^q$
\[
A(u_{pq})(\alpha*\beta)=\frac{(p+q)!}{p!\,q!}\alpha\square\beta
\quad\in A^{\mathfrak S_p\times\mathfrak S_q} \tag*{(29)}
\]
d'où, si on veut une relation dans $A^{\unsure{*|\bullet}}$ lui-même
\[
A(u_{p,p})(\alpha*\beta)=\frac{(\deleted{p+q}\inserted{2p})!}{p!\,q!}\alpha\beta \tag*{(30)}
\]
qui relie la multiplication dans les $A^n$ \deleted{\unsure{avec}}
à la multiplication ``globale'' $*$.

De plus, la formule 16 montre qu'on
a un hom canonique
\[
\left\{\begin{aligned}
&T\underset{\unsure{\cdot}}{=}T_G:A(G)\longrightarrow1+\widehat{A(G)^*}{}^+ \deleted{\unsure{z_+}}\\
&T(\alpha)=\sum T^n(\alpha)\deleted{\unsure{t^n}}\qquad(T^0(\alpha)=1)
\end{aligned}\right.\tag*{(31)}
\]
défini par
\[
T(\gamma(M))=\sum T^n(M)\deleted{\unsure{t^n}} \tag*{(32)}
\]
On a donc
\[
T^n(\alpha+\beta)=\sum_{p+q=n}T^p(\alpha)*T^q(\beta) \tag*{(33)}
\]
On déduit de (21)
\[
T^n(\alpha\beta)=T^n(\alpha)T^n(\beta) \tag*{(33)}
\]
(qui, compte tenu de (29), exprime ``presque''
% archive sheet number: 147
% printed form (vertical paper watermark, reverse face)
\unsure{GENTINE} \tikz[baseline=-.2ex,scale=.17]{\draw (0,0)--(1.6,0)--(1.8,.33)--(.5,.33)--(.32,.13)--(.92,.13);} TYPEWRITER \tikz[baseline=-.4ex,scale=.15]{\foreach \a in {0,45,...,315}{\draw (\a:.52) circle (.22);} \draw (0,0) circle (.21);} \unsure{ARGENT}
\scanpage{149}
$T^n(\alpha\beta)$ à l'aide de l'opération $*$ et des \deleted{\unsure{$\lambda$}} opérations
de restriction. Il faut voir comment s'expriment
les $T^n(\lambda^i(\alpha))$, on y reviendra.

Bien entendu, \deleted{\unsure{la} \illegible{1} \unsure{commutativité}} \inserted{les homomorphismes} $T_G^n$ sont fonctoriels :
si $u:G\longrightarrow G'$, on a \deleted{\unsure{cette}} commutativité \unsure{du|dans|des}
\diagram{\begin{tikzcd}[ampersand replacement=\&,column sep=large]
\text{(34)} \& A(G') \arrow[r,"T_{G'}"] \arrow[d,"A(u)"'] \& 1+\widehat{A(G')^*}{}^+ \arrow[d,"\widehat{A(u)}"]\\
\& A(G) \arrow[r,"T_G"] \& 1+\widehat{A(G)^*}{}^+
\end{tikzcd}}
% archive sheet number: 148
% printed form (vertical paper watermark, reverse face)
RITER \tikz[baseline=-.4ex,scale=.15]{\foreach \a in {0,45,...,315}{\draw (\a:.52) circle (.22);} \draw (0,0) circle (.21);} ARGENTINE \tikz[baseline=-.2ex,scale=.17]{\draw (0,0)--(1.6,0)--(1.8,.33)--(.5,.33)--(.32,.13)--(.92,.13);} \unsure{TYPEWRITE|TYPEWRITER}
\scanpage{150}
(i) Anneaux commutatifs [de Chern involutifs augmentés] $A^n$, $A^0=\mathbb Z$

(ii) Pour tout hom $u:\unsure{\mathfrak S|\mathcal S}_p\longrightarrow\mathfrak S_q$, un hom $A(u):A^q\longrightarrow A^p$
\[
A(uv)=A(v)A(u),\quad A(1_{\mathfrak S_p})=1_{A^p},\quad A(\sigma_g)=1_{A^p}\quad\text{si }g\in\mathfrak S_p.
\]

(iii) \deleted{Des} opérations bilinéaires $A^p\times A^q\longrightarrow A^{p+q}$, $(\alpha*\beta)$
faisant de $A^{\unsure{\bullet|*}}=\coprod_{\unsure{0}}^\infty A^n$ un anneau gradué commutatif
avec unité. On \unsure{:|a} la relation
\[
\frac{(2p)!}{(p!)^2}\alpha\beta=A(u_p)(\alpha*\beta)
\]
où $u_p$ est le composé naturel $\mathfrak S_p\xrightarrow{\Delta}\mathfrak S_p\times\mathfrak S_p\xrightarrow{u_{p,p}}\mathfrak S_{2p}$

(iv) \deleted{Des} homomorphisme de groupes
\[
T:A^1\longrightarrow\widetilde{A^*}=\mathbb Z\times(1+\widehat A^{*+})
\]
défini par des composantes $\varepsilon$ (l'augmentation de $A$)
et $T^n$, avec
\[
T^0(\alpha)=1,\quad T^1(\alpha)=\alpha,\quad T^n(\alpha)\in A^n
\]
On a donc
\[
T^n(\alpha+\beta)=\sum_{p+q=n}T^p(\alpha)T^q(\alpha)
\]
De plus
\[
T^n(\alpha.\beta)=T^n(\alpha)T^n(\beta)
\]

% Upside-down working, crossed through in diagonal strokes.
\diagram{\begin{tikzcd}[ampersand replacement=\&]
\deleted{G}\arrow[r,"\deleted{u}"]\arrow[d,"\deleted{\unsure{s|\alpha}}"'] \& \deleted{G'}\arrow[d,"\deleted{\beta}"]\\
\deleted{G}\arrow[r,"\deleted{\unsure{u'|u}}"] \& \deleted{G'}
\end{tikzcd}}
\[
\begin{aligned}
&\deleted{I^u[\mathcal A^{p_1}(M)\otimes\mathcal A^{p_2}(M')]}\\
&\deleted{I^{\unsure{u'}}(\mathcal A^{p_2}(M')\otimes\mathcal A^{p_1}(M))}\\
&\deleted{\unsure{M*M}}\\
&\deleted{\unsure{I}}\qquad\deleted{\unsure{M}}\qquad\deleted{\sigma_g:G\longrightarrow G}\\
&\deleted{T^n(\alpha)}\qquad\deleted{\unsure{M\otimes M'}}\qquad\deleted{M\longrightarrow I^{\unsure{\sigma_g}}M}\\
&\deleted{\unsure{\hat A}}\qquad\deleted{\lambda^n\alpha=\text{expression en les }T^n(\alpha)}
\end{aligned}
\]
% margin note written upside down
\margin{\uline{\deleted{\unsure{$A$}} $M'$ \unsure{sur}}}
% archive sheet number: 149
\scanpage{151}
\[
A(T)=\mathbb Z(\widehat T)
\]

Soit $\rho\unsure{\in|:}T\longrightarrow\operatorname{Gl}(n)$ / induit $\rho:T\longrightarrow k^{*\,n}$, d'où $n$ éléments
$\alpha_1,\ldots\alpha_n$ de $\widehat T$. Les fonctions symétriques élémentaires en
les $\alpha_i$ sont notées $C^i(\rho)$ $(0\leq i\leq n)$, ils ne dépendent
que de la classe de $\rho$. Soit $\unsure{C|\mathcal C}(\rho)=\sum_0^\infty C^i(\rho)\in\mathbb Z[\widehat T]$
($C^i(\unsure{\varphi|\rho})=0$ si $i>\deg\rho$) ; si $\rho$ est une extension de $\rho'$ par
$\rho''$, on a $C(\rho)=C(\rho')C(\rho'')$.

On peut prolonger $\rho\mapsto(\deg\rho,C(\rho))$ en un hom de groupe
\[
A(T)=\mathbb Z(\widehat T)\longrightarrow\mathbb Z\times(1+\mathbb Z[[\widehat T]]^+) \tag*{(1)}
\]
c'est un hom d'anneaux de Chern, caractérisé par
ses valeurs sur les $\rho_\chi$ $(\chi\in\widehat T)$
\[
\widetilde C(\rho_\chi)=(1,1+[\chi])
\]
Je dis que cet hom (1) est injectif \deleted{\unsure{(1)}} et à image
\deleted{\unsure{discr.}}dense. Pour le voir, on compose avec \deleted{\unsure{$\varphi$}}
\[
\varphi:\mathbb Z\times(1+\mathbb Z[[\widehat T]]^+)\longrightarrow\mathbb Q[[\widehat T]]
\]
transformant $\chi$ en $e^\chi$, et on note que les
$e^\chi$ sont linéairement indépendants. Pour montrer que
l'image \deleted{est} dense, on montre qu'on peut trouver
un hom $\rho$ dont les $C^1,\ldots C^n\unsure{\mid|)}$ soient donnés,
i.e \deleted{des formes} étant donné des \deleted{formes sur $\mathbb Z^r$}
polynomes $C^i\in\mathbb Z[X_1,\ldots,X_r]$ $(1\leq i\leq n)$, $C^i$
de degré $i$, on peut trouver des formes linéaires
$\alpha^1,\ldots,\alpha^N$ en les $X_i$, telles que
$C^i=\sigma_i(\alpha^1,\ldots,\alpha^N)$ pour $1\leq i\leq n$. Raisonnons sur
% archive sheet number: 150
% printed form (vertical paper watermark, reverse face)
\tikz[baseline=-.4ex,scale=.15]{\foreach \a in {0,45,...,315}{\draw (\a:.52) circle (.22);} \draw (0,0) circle (.21);} ARGENTINE \tikz[baseline=-.2ex,scale=.17]{\draw (0,0)--(1.6,0)--(1.8,.33)--(.5,.33)--(.32,.13)--(.92,.13);} TYPEWRITER \tikz[baseline=-.4ex,scale=.15]{\foreach \a in {0,45,...,315}{\draw (\a:.52) circle (.22);} \draw (0,0) circle (.21);} \unsure{AR}
\scanpage{152}
$n$, trivial pour $n=0$. Cela nous ramène au cas
où $C^1=\ldots=C^{n-1}=0$, donc à prouver le

\uline{Lemme} Il existe des formes linéaires $\alpha^1,\ldots,\alpha^N$
entières en $X^1,\ldots X^r$ telles que
\[
\sigma^1((\alpha^i))=\ldots=\sigma^{n-1}((\unsure{\alpha_i|\alpha^i}))=0,\quad
\sigma^n((\unsure{\alpha_i|\alpha^i}))=C^n\ \text{polynome donné}
\]

\diagram{\begin{tikzcd}[ampersand replacement=\&,row sep=1.6em,column sep=large,execute at end picture={\draw (current bounding box.north west) -- (current bounding box.north east) -- (current bounding box.south east);}]
\mathrm Z[\unsure{\xi_1,\ldots,\xi_n|\zeta_1,\ldots,\zeta_n|\xi_1,\ldots}] \& {}\\
E_T\arrow[d] \& {}\\
B_T \& {}\\
\mathbb Z[\widehat T]\arrow[r,"\unsure{\sim}"] \& H^*(B_T,\mathbb Z)\\
{} \& A(T)\arrow[ul,"\text{alg}"]\arrow[u,"\text{top via Chern}"']\\
\mathbb Z[\widehat T]\arrow[r,"{\text{\unsure{hom. car.}}}"{draw,circle,inner sep=1pt,above}] \& H^*(X,\mathbb Z)\\
{} \& A(T)\arrow[ul]\arrow[u,"\text{top. via Chern}"']
\end{tikzcd}}

Il faudrait prouver que
$\mathbb Z[\widehat T]$ est l'anneau de
Chern attaché à $A(T)$
$\simeq\mathbb Z(\widehat T)$
% archive sheet number: 151
% printed form (vertical paper watermark)
\tikz[baseline=-.4ex,scale=.15]{\foreach \a in {0,45,...,315}{\draw (\a:.52) circle (.22);} \draw (0,0) circle (.21);} ARGENTINE \tikz[baseline=-.2ex,scale=.17]{\draw (0,0)--(1.6,0)--(1.8,.33)--(.5,.33)--(.32,.13)--(.92,.13);} TYPEWRITER \tikz[baseline=-.4ex,scale=.15]{\foreach \a in {0,45,...,315}{\draw (\a:.52) circle (.22);} \draw (0,0) circle (.21);} \unsure{AR}
\scanpage{153}
Soit $G$ un groupe algébrique. On désigne par $\deleted{A}K(G)$
l'anneau \deleted{\unsure{des}} engendré par les classes des représentations
rationnelles irréductibles de $G$ : en tant que groupe
additif, il est donc libre ; c'est un \uline{anneau}
\uline{de Chern}. On désigne par $C\deleted{A}K(G)$ ``l'anneau
des classes de Chern'' de $\deleted{A}K(G)$, \inserted{c'est un anneau gradué,} \deleted{$\unsure{A(G)}$ et $\unsure{CA(G)}$ sont des foncteurs}
et on a un homomorphisme canonique
\[
\deleted{A}K(G)\xrightarrow{c_G}\mathbb Z\times(1+C\deleted{A}K(G)^+) \tag*{(1)}
\]
$\deleted{A}K(G)$ et $C\deleted{A}K(G)$ sont des foncteurs en $G$, et
(1) est un hom. fonctoriel. Si $\sigma$ est un
autom. \uline{intérieur} de $G$, alors $\sigma^*$ est l'automorphisme identique de $K(G)$ et $CK(G)$.

Soit $P$ un fibré algébrique principal de
groupe structural $G$, de base $X$. À toute
\deleted{\unsure{quel}} représentation lin. rationnelle $\rho$ de $G$
est associé un fibré vectoriel \inserted{algébrique} sur $X$, d'où
un élément de $\deleted{\unsure{C}}\deleted{A}K(X)\deleted{\unsure{{}_P}}$. On en
conclut un hom. d'anneaux de Chern
\[
K(G)\longrightarrow K(X)
\]
d'où un hom d'anneaux gradués
\[
CK(G)\longrightarrow CK(X)=A(X) \tag*{(2)}
\]
appelé \uline{hom. caractéristique} du fibré $P$. On
a commutativité des
% archive sheet number: 152
% printed form (vertical paper watermark)
\unsure{R|ER} \tikz[baseline=-.4ex,scale=.15]{\foreach \a in {0,45,...,315}{\draw (\a:.52) circle (.22);} \draw (0,0) circle (.21);} ARGENTINE \tikz[baseline=-.2ex,scale=.17]{\draw (0,0)--(1.6,0)--(1.8,.33)--(.5,.33)--(.32,.13)--(.92,.13);} TYPEWRITER \tikz[baseline=-.4ex,scale=.15]{\foreach \a in {0,45,...,315}{\draw (\a:.52) circle (.22);} \draw (0,0) circle (.21);} \unsure{A}
\scanpage{154}
\diagram{\begin{tikzcd}[ampersand replacement=\&,column sep=large]
\text{(3)} \& K(G)\arrow[r,"\text{ass.}"]\arrow[d,"c_G"] \& K(X)\arrow[d,"c_X"]\\
\& \widetilde{CK}(G)\arrow[r,"\text{hom. car.}","\text{pour }P"'] \& \widetilde{CK}(X)=\widetilde A(X)
\end{tikzcd}}

De plus, si on a un hom. rat. $G\xrightarrow{\unsure{\rho|f|\varphi}}G'$
et si $P'$ est associé à $P$ et $\rho$, on a commutativité dans
\diagram{\begin{tikzcd}[ampersand replacement=\&]
\text{(4)} \& CK(G')\arrow[rr,"\rho^*"]\arrow[dr,"\unsure{\gamma|r|\tau}_{P'}"] \& {} \& C\deleted{\unsure{A}}K(G)\arrow[dl,"\gamma_P"']\\
\& {} \& A(X) \& {}
\end{tikzcd}}

Il reste à calculer explicitement $CK(G)$.

Voici les résultats essentiels dans cette voie.

\uline{Proposition} Soit $T$ un tore algébrique \inserted{de dim. $r$}. Alors
\[
\left|\begin{aligned}
&K(T)\simeq\mathbb Z(\widehat T)\qquad(\text{algèbre du groupe }\widehat T)\\
&\big\{\simeq\mathbb Z[x^1,\ldots,x^r]_{x^1\cdots x^r}\ \deleted{\text{avec}}\inserted{\text{avec}}\quad
\varepsilon(x^i)=1,\ \lambda x^i=1-x^i t\big\}\\
&\text{avec }\quad\varepsilon(\chi)=1,\quad\lambda(\chi)=1+\chi t\quad\text{pour }\chi\in\widehat T\\
&\text{On a }\quad CK(T)\simeq\mathbb Z[\widehat T]\qquad(\text{algèbre symétrique}\\
&\text{sur }\mathbb Z\text{ du }\mathbb Z\text{-module }\widehat T),\quad\text{avec}\\
&\qquad c_T(\chi)=1+\chi\qquad\text{i.e. }c_T^i(\chi)=
\begin{cases}1&\text{si }i=0\\\chi&\text{si }i=1\\0&\text{si }i>\unsure{1|0}\end{cases}
\end{aligned}\right.
\]
\begin{center}-- -- -- -- -- -- -- --\end{center}
% archive sheet number: 153
% printed form (vertical paper watermark, reverse face)
\unsure{R|ER} \tikz[baseline=-.4ex,scale=.15]{\foreach \a in {0,45,...,315}{\draw (\a:.52) circle (.22);} \draw (0,0) circle (.21);} ARGENTINE \tikz[baseline=-.2ex,scale=.17]{\draw (0,0)--(1.6,0)--(1.8,.33)--(.5,.33)--(.32,.13)--(.92,.13);} TYPEWRITER \tikz[baseline=-.4ex,scale=.15]{\foreach \a in {0,45,...,315}{\draw (\a:.52) circle (.22);} \draw (0,0) circle (.21);} \unsure{A}
